\subsection{Solvability, equivariantization and module ranks}
We give three further general exclusions using only small groups and exact fusion rules. By \cite[Theorem~1.6]{ENOSolvable}, a fusion category of FP dimension $p^a q^b$ is solvable. The characterization in its Definition~1.2 and Proposition~4.4 expresses it as successive cyclic-prime extensions and equivariantizations. If its adjoint subcategory is the whole category, its last nontrivial step cannot be an extension. Thus it is an actual $C_p$-equivariantization, with a central $\Rep(C_p)$ whose underlying invertibles embed into the category. We use the orbit/stabilizer description of equivariant simples \cite[Section~2]{BurciuNatale}; cyclic stabilizers have only one-dimensional irreducible projective representations.

\begin{proposition}\label{prop:dim24}
The ring $R8\_1\_29$ has no complex fusion categorification.
\end{proposition}
\begin{proof}
Its invertibles form $C_6$, its other two simples $X,Y$ are self-dual of dimension three, and its adjoint ring is the whole ring. Its dimension is $24$, so a putative category $\mathcal C$ is $\mathcal D^{C_p}$ with $p=2$ or $3$.

For $p=3$, the order-three character subgroup fixes $X,Y$ and has two free orbits on the six invertibles. The forgetful images of $X,Y$ split into three distinct invertibles each; those of the invertibles give two more. Thus $\mathcal D=\Vec_A^\omega$, $|A|=8$, with two free $C_3$-orbits. The only groups of order eight admitting an order-three automorphism are $C_2^3$ and $Q_8$. On $C_2^3$ the action has a fixed line and an irreducible plane. The three nonzero elements in that plane form a free orbit whose product has coefficient two into the same free orbit. The output stabilizer is trivial, so associator and tensorator choices cannot change this coefficient. On $Q_8$, inversion exchanges the orbits $\{i,j,k\}$ and $\{-i,-j,-k\}$; their two induced dimension-three objects are distinct duals. These conclusions contradict multiplicity one and self-duality, respectively.

For $p=2$, the character involution exchanges $X,Y$. De-equivariantization has the fusion rules
\[
G(\mathcal D)=C_3=\langle a\rangle,\qquad a\rho=\rho a=\rho,
\qquad \rho^2=1+a+a^2+2\rho.
\]
It has dimension twelve and trivial universal grading. The same last-step argument gives
$\mathcal D=(\Vec_V^\omega)^{C_3}$, where $V=C_2^2$ and $C_3$ permutes the three nonzero elements transitively. No assumption on $\omega$ is made. A monoidal automorphism of $\mathrm{Id}_{\mathcal D}$ has scalar one on $\rho$, because $\rho$ occurs in $\rho^2$, and then on all invertibles. Consequently the map from central invertibles to underlying invertibles is injective, and the central $\Rep(C_3)$ is unique.

The $C_2$ action on $\mathcal D$ fixes every simple: exchanging $a,a^2$ would produce a dimension-two simple upstairs. Every such tensor autoequivalence preserves the unique central $\Rep(C_3)$. Its restriction is tensor-isomorphic to the identity, since $H^2(C_3,\C^\times)=0$. Equivariantization/de-equivariantization \cite[Theorem~4.18(i) and Proposition~4.19]{DGNO} therefore lifts it to a $C_3$-equivariant tensor autoequivalence of $\Vec_V^\omega$. Its object permutation centralizes $C_3$ in $\operatorname{Aut}(V)=S_3$, so lies in $C_3$. Composing with the inverse action functor removes that permutation without changing the descended equivalence class.

Identity-on-objects tensor autoequivalences of $\Vec_V^\omega$ are classified by $H^2(V,\C^\times)=C_2$: the two associators cancel in the tensorator equation. Their equivariant structures have ambiguity $H^1(C_3,\widehat V)=0$, since multiplication by three is invertible on $\widehat V$. Both possible classes have explicit representatives. Besides the identity, take
\[
J(u,v)=(-1)^{u_1v_1+u_1v_2+u_2v_2},\qquad u,v\in\mathbb F_2^2.
\]
This nontrivial normalized $2$-cocycle is $C_3$-invariant and satisfies $J^2=1$. Thus it commutes even with nontrivial scalar tensorators of the existing $C_3$ action and descends with identity comparisons. There are no further $C_2$ coherence choices, because $\operatorname{Aut}_{\otimes}(\mathrm{Id}_{\mathcal D})=1$.

The identity action leaves the coefficient two in $\rho^2$. For the $J$ action, all degree-zero channels have $J(u,u)=-1$ for $u\ne0$; for each nonzero output degree, the two ordered channels have opposite signs. Hence each lift $\rho_\pm$ has square containing the three negative-character invertibles and one copy of each $\rho_\pm$. The unit occurs in $\rho_+\rho_-$, not in either square. The two lifts are distinct duals, again contradicting the self-duality in $R8\_1\_29$. This exhausts both prime branches for arbitrary pointed associators.
\end{proof}

\begin{lemma}[Module-rank divisibility]\label{lem:modulegcd}
For a fusion category $\mathcal B$, let $g_{\mathcal B}$ be the greatest common divisor of the ranks of its indecomposable finite semisimple module categories. If $\mathcal B\subseteq\mathcal D$ is a fusion subcategory, then $g_{\mathcal B}\mid\rank(\mathcal D)$.
\end{lemma}
\begin{proof}
Decompose $\mathcal D$, with its left tensor action by $\mathcal B$, into indecomposable module categories. Their ranks add to $\rank(\mathcal D)$. This statement uses ranks in the specified module presentation; it makes no assertion that module ranks are Morita invariant.
\end{proof}

\begin{proposition}\label{prop:dim72}
The rings $R8\_1\_45$ and $R8\_1\_47$ have no complex fusion categorification.
\end{proposition}
\begin{proof}
Each ring has simple dimensions $(1,1,2,3,3,4,4,4)$ and admits labels $1,h,E,P,Q,W_0,W_1,W_2$ with
\[
\begin{gathered}
h^2=1,\quad hE=E,\quad hP=Q,\quad hW_i=W_i,\quad P^*=Q,\\
E^2=1+h+E,\qquad EP=EQ=P+Q,\qquad PQ=1+E+P+Q.
\end{gathered}
\]
The adjoint ring is the whole ring. Dimension $72$ and the invertible group $C_2$ therefore force a putative categorification to be $\mathcal C=\mathcal D^{C_2}$ with $\FPdim(\mathcal D)=36$. The forgetful images are
\[
F(1)=F(h)=1,\quad F(E)=a+a^2,\quad F(P)=F(Q)=Y,
\quad F(W_i)=U_i+V_i.
\]
Here $a$ has order three, $Y$ is simple of dimension three, and the six $U_i,V_i$ are distinct simples of dimension two. Thus $\rank(\mathcal D)=10$. The full fusion subcategory $\mathcal B$ on $1,a,a^2,Y$ satisfies
\[
aY=Ya=Y,\qquad Y^*=Y,\qquad Y^2=1+a+a^2+2Y.
\]
The $C_2$ action exchanges $a,a^2$ and fixes $Y$; its equivariantization inside $\mathcal C$ has the displayed nonselfdual $P,Q$.

As in the preceding proof, solvability and trivial universal grading imply
$\mathcal B=(\Vec_V^\omega)^{C_3}$, $V=C_2^2$, with the transitive action. The restriction map
\[
H^3(V,\C^\times)\longrightarrow
\prod_{\substack{L\le V\\|L|=2}}H^3(L,\C^\times)
\]
is an isomorphism $C_2^3\to C_2^3$. In integral-cohomology generators $u^2,v^2,uv$, the three restrictions of $(r,s,t)$ are $(r,s,r+s+t)$. Since $C_3$ permutes the factors, there are exactly two invariant classes: zero and the class nontrivial on every subgroup of order two.

Suppose first $[\omega]=0$. The lift of the object $C_3$ action is unique in $\operatorname{Aut}_{\otimes}(\Vec_V)=C_2\times S_3$, and its coherence is unique because $H^2(C_3,\widehat V)=0$. Therefore $\mathcal B\simeq\Rep(A_4)$. Its central $\Rep(C_3)$ is unique, since $\operatorname{Aut}_{\otimes}(\mathrm{Id}_{\mathcal B})=1$, and is preserved by the $C_2$ action. Sequential equivariantization and de-equivariantization \cite[Section~2.6 and Proposition~2.10]{ENOSolvable}\cite[Proposition~4.19]{DGNO} identify $\mathcal B^{C_2}$ with $(\Vec_V)^{S_3}$. The action on nonzero elements is the usual faithful $S_3$ permutation action; inversion on the character $C_3$ forces the group $C_3\rtimes C_2=S_3$.

There are exactly two lifts of this object action to $C_2\times S_3$, given by the trivial or parity homomorphism into $H^2(V,\C^\times)=C_2$. Their coherence ambiguity vanishes: $H^j(S_3,\widehat V)=0$ for $j>0$, by the normal $C_3$ subgroup, whose positive cohomology and fixed space on $\widehat V$ both vanish. Both lifts yield self-dual dimension-three orbit objects. For completeness choose a stabilizing transposition $\tau$ swapping coordinates and its fixed nonzero element $v=(1,1)$. For the trivial lift, the equivariance scalars $u=\pm1$ act trivially on the unit channel of $v\otimes v$. For the other lift choose tensorator $\mu(a,b)=(-1)^{a_1b_2}$. Its square is trivialized by $\eta_a=(-1)^{a_1a_2}$, so $\eta_v=\mu(v,v)=-1$. The equivariance scalar satisfies $u^2=\eta_v$ and acts on the unit channel by $u^2/\mu(v,v)=1$. A $\tau$-invariant character change of $\eta$ is one on $v$, so cannot change the conclusion. Frobenius reciprocity for the orbit induction gives a tensor unit in each of the two squares. This contradicts $P^*=Q$.

It remains that $[\omega]$ is nontrivial on every order-two subgroup. The crossed product of $\Vec_V^\omega$ by $C_3$ is pointed, and realizes $\mathcal B$ as the group-theoretical category
\[
\mathcal B\simeq\mathcal C(A_4,\Omega,K,\psi),\qquad K=C_3,
\qquad [\Omega]|_V=[\omega],\quad [\Omega]|_K=0.
\]
By the classification of pointed module categories and its double-coset rank formula \cite[Example~2.1, Theorem~3.1 and Proposition~3.1]{OstrikDouble}, indecomposable $\mathcal B$-module categories correspond to subgroups $H\le A_4$ with $[\Omega]|_H=0$ and permissible cochains. If $|H|$ were even, Cauchy's theorem would give a subgroup of order two, contradicting the restriction of $\Omega$. Thus $H$ is trivial or has order three. For $H=1$, the four double cosets in $K\backslash A_4/H$ have trivial intersections, giving module rank four. For $H=C_3$, conjugate it to $K$. There are two double cosets, with intersections $C_3$ and $1$; projective cocycles on $C_3$ are coboundaries, so their contributions are three and one. Every indecomposable $\mathcal B$-module category therefore has rank four, regardless of the permissible cochains. Lemma~\ref{lem:modulegcd} would force $4\mid\rank(\mathcal D)=10$, a contradiction.
\end{proof}

The portable verifiers check the source-bound tensors, the displayed fusion identities, the complete small cocycle identities, and all relevant $A_4$ subgroups and double cosets. These finite checks certify the hypotheses; the categorical exhaustion is supplied by the proofs above. Neither argument assumes a braiding or a pivotal structure.
